\documentclass[10pt,a4paper]{amsart}
\usepackage{amsmath,amssymb,amsthm,graphicx}
% For commutative diagrams we use Xy-pic
\usepackage[all]{xy}

% We use 2cell for 2-commutative diagrams.
\xyoption{2cell}
% Theorem environments.
%
\theoremstyle{plain}
\newtheorem{theorem}[subsection]{Theorem}
\newtheorem{proposition}[subsection]{Proposition}
\newtheorem{lemma}[subsection]{Lemma}

\theoremstyle{definition}
\newtheorem{definition}[subsection]{Definition}
\newtheorem{example}[subsection]{Example}
\newtheorem{exercise}[subsection]{Exercise}
\newtheorem{situation}[subsection]{Situation}

\theoremstyle{remark}
\newtheorem{remark}[subsection]{Remark}
\newtheorem{remarks}[subsection]{Remarks}

\numberwithin{equation}{subsection}

% Macros
%
\def\lim{\mathop{\mathrm{lim}}\nolimits}
\def\colim{\mathop{\mathrm{colim}}\nolimits}
\def\Spec{\mathop{\mathrm{Spec}}}
\def\Hom{\mathop{\mathrm{Hom}}\nolimits}
\def\Ext{\mathop{\mathrm{Ext}}\nolimits}
\def\SheafHom{\mathop{\mathcal{H}\!\mathit{om}}\nolimits}
\def\SheafExt{\mathop{\mathcal{E}\!\mathit{xt}}\nolimits}
\def\Sch{\mathit{Sch}}
\def\Mor{\mathop{\mathrm{Mor}}\nolimits}
\def\Ob{\mathop{\mathrm{Ob}}\nolimits}
\def\Sh{\mathop{\mathit{Sh}}\nolimits}
\def\NL{\mathop{N\!L}\nolimits}
\def\CH{\mathop{\mathrm{CH}}\nolimits}
\def\proetale{{pro\text{-}\acute{e}tale}}
\def\etale{{\acute{e}tale}}
\def\QCoh{\mathit{QCoh}}
\def\Ker{\mathop{\mathrm{Ker}}}
\def\Im{\mathop{\mathrm{Im}}}
\def\Coker{\mathop{\mathrm{Coker}}}
\def\Coim{\mathop{\mathrm{Coim}}}

% Boxtimes
%
\DeclareMathSymbol{\boxtimes}{\mathbin}{AMSa}{"02}

%
% Macros for moduli stacks/spaces
%
\def\QCohstack{\mathcal{QC}\!\mathit{oh}}
\def\Cohstack{\mathcal{C}\!\mathit{oh}}
\def\Spacesstack{\mathcal{S}\!\mathit{paces}}
\def\Quotfunctor{\mathrm{Quot}}
\def\Hilbfunctor{\mathrm{Hilb}}
\def\Curvesstack{\mathcal{C}\!\mathit{urves}}
\def\Polarizedstack{\mathcal{P}\!\mathit{olarized}}
\def\Complexesstack{\mathcal{C}\!\mathit{omplexes}}
% \Pic is the operator that assigns to X its picard group, usage \Pic(X)
% \Picardstack_{X/B} denotes the Picard stack of X over B
% \Picardfunctor_{X/B} denotes the Picard functor of X over B
\def\Pic{\mathop{\mathrm{Pic}}\nolimits}
\def\Picardstack{\mathcal{P}\!\mathit{ic}}
\def\Picardfunctor{\mathrm{Pic}}
\def\Deformationcategory{\mathcal{D}\!\mathit{ef}}
\newenvironment{reference}{}{}
\newenvironment{slogan}{}{}
\newenvironment{history}{}{}
\title{The Picard group of a numerical type embeds into the cokernel of its intersection matrix}
\author{The Stacks Project Authors}
\date{stacks-project:models.tex:local-snapshot-20261012; GNU FDL 1.2 or later}
\begin{document}
\maketitle
\noindent\textit{Editorial note.} The counted unit is the Stacks Project Lemma lemma-picard-T-and-A at source lines 1143-1147 together with its proof at source lines 1149-1166, a straight adjacent pairing in the chapter file. The chapter is the upstream file models.tex, which is the file name recorded in the source record; the pool directory carries the longer name stacks-models and is not used as the chapter name anywhere in this card.\par\medskip
\noindent\textit{Editorial note.} This chapter's candidate index records the statement body and the proof body without their closing markers, so the retained spans used here are the full environments: the statement span ends at the line of the closing marker of the lemma environment and the proof span starts at the opening marker of the proof environment. The same one-line accounting was recorded in the earlier cards of this chapter.\par\medskip
\noindent\textit{Editorial note.} The retained proof draws a commutative diagram with two exact rows using the author's Xy-pic matrix syntax, so the retained preamble keeps the chapter preamble's Xy-pic block at lines 14-18 verbatim, including the option that enables two-cell diagrams. Nothing about the diagram has been rewritten or replaced.\par\medskip
\noindent\textit{Editorial note.} Neither the retained statement nor the retained proof contains a reference to another labelled result of the chapter and neither contains a citation, so this card carries no navigation pointer and no bibliography entry.\par\medskip
\noindent\textit{Editorial note.} The retained preamble is the chapter preamble's Xy-pic setup at lines 14-18, the theorem declarations and mathematical macros at lines 154-217, and the three no-op wrapper environments.\par\medskip
\noindent\textit{Editorial note.} The proof is the author's own complete argument and contains no omission wording: the presentation of the Picard group, the two exact rows and the application of the snake lemma are all written out. The author's own wording attributes the conclusion directly to the snake lemma, which is a standard tool and is not a missing step.\par\medskip
\noindent\textit{Allowed standard prerequisites (editorial).}\begin{itemize}
\item Numerical types and their Picard group. A numerical type is a collection of integers n, m of subscript i, a of subscript i j, w of subscript i and g of subscript i satisfying the symmetry and recurrence conditions of the chapter, and its Picard group is defined as the cokernel of the matrix of the rational quotients a of subscript i j over w of subscript i; the author's own notation is retained inside the counted statement and proof.. The counted lemma compares that presentation of the Picard group with the cokernel of the plain intersection matrix, and the comparison is the input for the rank computation and the embedding statements that follow in the same section.
\end{itemize}\medskip
\begin{lemma}
\label{lemma-picard-T-and-A}
Let $n, m_i, a_{ij}, w_i, g_i$ be a numerical type $T$.
Then $\Pic(T) \subset \Coker(A)$ where $A = (a_{ij})$.
\end{lemma}

\begin{proof}
Since $\Pic(T)$ is the cokernel of $(a_{ij}/w_i)$
we see that there is a commutative diagram
$$
\xymatrix{
0 \ar[r] &
\mathbf{Z}^{\oplus n} \ar[rr]_A & &
\mathbf{Z}^{\oplus n} \ar[rr] & &
\Coker(A) \ar[r] & 0 \\
0 \ar[r] &
\mathbf{Z}^{\oplus n} \ar[rr]^{(a_{ij}/w_i)} \ar[u]_{\text{id}} & &
\mathbf{Z}^{\oplus n} \ar[rr] \ar[u]_{\text{diag}(w_1, \ldots, w_n)} & &
\Pic(T) \ar[r] \ar[u] & 0
}
$$
with exact rows. By the snake lemma we conclude that
$\Pic(T) \subset \Coker(A)$.
\end{proof}

\end{document}
